% Wspólna bibliografia Basilisk Systems --- % Wspólna bibliografia Basilisk Systems --- % Wspólna bibliografia Basilisk Systems --- % Wspólna bibliografia Basilisk Systems --- \input{basilisk_bib} w podręczniku decyzji i komentarzu do umowy spółki.
% Odesłanie w tekście: \Zr{klucz} -> [ETYKIETA] z linkiem do pozycji.
\makeatletter
\providecommand{\BibItem}[2]{\item[\textbf{[#1]}]\def\@currentlabel{#1}\def\@currentlabelname{#1}\phantomsection\label{bib:#2}}
\makeatother
\providecommand{\Zr}[1]{[\ref{bib:#1}]}
\newenvironment{BibList}{\begin{list}{}{\setlength{\labelwidth}{2.6cm}\setlength{\leftmargin}{2.9cm}\setlength{\labelsep}{0.3cm}\setlength{\itemsep}{0.15em}\setlength{\parsep}{0pt}\small}}{\end{list}}

\chapter*{Bibliografia}
\addcontentsline{toc}{chapter}{Bibliografia}

Akty prawne podano według miejsca pierwotnej publikacji oraz ostatniego tekstu jednolitego (stan: listopad 2026), z dopiskiem „ze zm.”; przed użyciem należy sprawdzić, czy nie ogłoszono nowszego tekstu w Internetowym Systemie Aktów Prawnych (isap.sejm.gov.pl) albo w Dzienniku Ustaw (eli.gov.pl). Tezy orzeczeń podano w skrócie; przed powołaniem się na nie w sporze należy sprawdzić pełny tekst.

\section*{A. Akty prawne krajowe}
\begin{BibList}
\BibItem{KSH}{ksh} Ustawa z dnia 15 września 2000 r. --- Kodeks spółek handlowych (Dz.U. 2000 nr 94 poz. 1037 ze zm.; t.j. Dz.U. 2024 poz. 18 ze zm., w tym ustawą z dnia 23 stycznia 2026 r., Dz.U. 2026 poz. 176 --- zmiany rejestru akcjonariuszy od 18 lutego 2027 r.); przepisy o prostej spółce akcyjnej: dział IA, art.~300$^{1}$--300$^{134}$.
\BibItem{Nowela PSA}{psa2019} Ustawa z dnia 19 lipca 2019 r. o zmianie ustawy --- Kodeks spółek handlowych oraz niektórych innych ustaw (Dz.U. 2019 poz. 1655) --- wprowadzenie prostej spółki akcyjnej; ustawa w pierwotnym brzmieniu wchodziła w życie 1 marca 2020 r., a przepisy o PSA stosuje się od 1 lipca 2021 r. (po przesunięciu terminu).
\BibItem{KC}{kc} Ustawa z dnia 23 kwietnia 1964 r. --- Kodeks cywilny (Dz.U. 1964 nr 16 poz. 93; t.j. Dz.U. 2026 poz. 795 ze zm.).
\BibItem{KRO}{kro} Ustawa z dnia 25 lutego 1964 r. --- Kodeks rodzinny i opiekuńczy (Dz.U. 1964 nr 9 poz. 59; t.j. Dz.U. 2026 poz. 236 ze zm.).
\BibItem{KPC}{kpc} Ustawa z dnia 17 listopada 1964 r. --- Kodeks postępowania cywilnego (Dz.U. 1964 nr 43 poz. 296; t.j. Dz.U. 2026 poz. 468 ze zm.).
\BibItem{KP}{kp} Ustawa z dnia 26 czerwca 1974 r. --- Kodeks pracy (Dz.U. 1974 nr 24 poz. 141; t.j. Dz.U. 2025 poz. 277 ze zm.).
\BibItem{PrAut}{praut} Ustawa z dnia 4 lutego 1994 r. o prawie autorskim i prawach pokrewnych (Dz.U. 1994 nr 24 poz. 83; t.j. Dz.U. 2025 poz. 24 ze zm.).
\BibItem{PWP}{pwp} Ustawa z dnia 30 czerwca 2000 r. --- Prawo własności przemysłowej (Dz.U. 2001 nr 49 poz. 508; t.j. Dz.U. 2023 poz. 1170 ze zm.).
\BibItem{UZNK}{uznk} Ustawa z dnia 16 kwietnia 1993 r. o zwalczaniu nieuczciwej konkurencji (Dz.U. 1993 nr 47 poz. 211; t.j. Dz.U. 2026 poz. 85 ze zm.).
\BibItem{PIT}{pit} Ustawa z dnia 26 lipca 1991 r. o podatku dochodowym od osób fizycznych (Dz.U. 1991 nr 80 poz. 350; t.j. Dz.U. 2026 poz. 592 ze zm.).
\BibItem{CIT}{cit} Ustawa z dnia 15 lutego 1992 r. o podatku dochodowym od osób prawnych (Dz.U. 1992 nr 21 poz. 86 ze zm.).
\BibItem{UoR}{uor} Ustawa z dnia 29 września 1994 r. o rachunkowości (Dz.U. 1994 nr 121 poz. 591; t.j. Dz.U. 2026 poz. 522 ze zm.).
\BibItem{UKRS}{ukrs} Ustawa z dnia 20 sierpnia 1997 r. o Krajowym Rejestrze Sądowym (Dz.U. 1997 nr 121 poz. 769; t.j. Dz.U. 2025 poz. 869 ze zm.).
\BibItem{UONI}{uoni} Ustawa z dnia 5 sierpnia 2010 r. o ochronie informacji niejawnych (Dz.U. 2010 nr 182 poz. 1228; t.j. Dz.U. 2025 poz. 1209 ze zm.).
\BibItem{UOTS}{uots} Ustawa z dnia 29 listopada 2000 r. o obrocie z zagranicą towarami, technologiami i usługami o znaczeniu strategicznym dla bezpieczeństwa państwa, a także dla utrzymania międzynarodowego pokoju i bezpieczeństwa (Dz.U. 2000 nr 119 poz. 1250; t.j. Dz.U. 2023 poz. 1582 ze zm., w tym Dz.U. 2026 poz. 471 --- definicje i zezwolenia, m.in. pomoc techniczna).
\BibItem{UWT}{uwt} Ustawa z dnia 13 czerwca 2019 r. o wykonywaniu działalności gospodarczej w zakresie wytwarzania i obrotu materiałami wybuchowymi, bronią, amunicją oraz wyrobami i technologią o przeznaczeniu wojskowym lub policyjnym (Dz.U. 2019 poz. 1214; t.j. Dz.U. 2023 poz. 1743 ze zm.); koncesja na obrót technologią o przeznaczeniu wojskowym (art.~7), wymogi wobec wspólników z co najmniej 20\% akcji (art.~10 i 17)..
\BibItem{UKI}{uki} Ustawa z dnia 24 lipca 2015 r. o kontroli niektórych inwestycji (Dz.U. 2015 poz. 1272; t.j. Dz.U. 2026 poz. 47 ze zm.).
\BibItem{USank}{usank} Ustawa z dnia 13 kwietnia 2022 r. o szczególnych rozwiązaniach w zakresie przeciwdziałania wspieraniu agresji na Ukrainę oraz służących ochronie bezpieczeństwa narodowego (Dz.U. 2022 poz. 835; t.j. Dz.U. 2025 poz. 514 ze zm.).
\BibItem{USyg}{usyg} Ustawa z dnia 14 czerwca 2024 r. o ochronie sygnalistów (Dz.U. 2024 poz. 928).
\BibItem{UODO}{uodo} Ustawa z dnia 10 maja 2018 r. o ochronie danych osobowych (Dz.U. 2018 poz. 1000 ze zm.).
\BibItem{PrLot}{prlot} Ustawa z dnia 3 lipca 2002 r. --- Prawo lotnicze (Dz.U. 2002 nr 130 poz. 1112; t.j. Dz.U. 2025 poz. 1431 ze zm.).
\BibItem{PKD 2025}{pkd} Rozporządzenie Rady Ministrów z dnia 18 grudnia 2024 r. w sprawie Polskiej Klasyfikacji Działalności (PKD) (Dz.U. 2024 poz. 1936).
\BibItem{WT}{wt} Rozporządzenie Rady Ministrów z dnia 17 września 2019 r. w sprawie klasyfikacji rodzajów materiałów wybuchowych, broni, amunicji oraz wyrobów i technologii o przeznaczeniu wojskowym lub policyjnym, na których wytwarzanie lub obrót jest wymagane uzyskanie koncesji (Dz.U. 2019 poz. 1888) --- obejmuje m.in. bezzałogowe statki powietrzne i ich systemy (WT~V pkt~3).
\end{BibList}

\section*{B. Akty prawne Unii Europejskiej}
\begin{BibList}
\BibItem{RODO}{rodo} Rozporządzenie Parlamentu Europejskiego i Rady (UE) 2016/679 z dnia 27 kwietnia 2016 r. w sprawie ochrony osób fizycznych w związku z przetwarzaniem danych osobowych i w sprawie swobodnego przepływu takich danych oraz uchylenia dyrektywy 95/46/WE (ogólne rozporządzenie o ochronie danych) (Dz.Urz. UE L 119 z 4.5.2016, s.~1).
\BibItem{Dual-use}{dualuse} Rozporządzenie Parlamentu Europejskiego i Rady (UE) 2021/821 z dnia 20 maja 2021 r. ustanawiające unijny system kontroli wywozu, pośrednictwa, pomocy technicznej, tranzytu i transferu produktów podwójnego zastosowania (wersja przekształcona) (Dz.Urz. UE L 206 z 11.6.2021, s.~1, ze zm.).
\BibItem{EDF}{edf} Rozporządzenie Parlamentu Europejskiego i Rady (UE) 2021/697 z dnia 29 kwietnia 2021 r. ustanawiające Europejski Fundusz Obronny i uchylające rozporządzenie (UE) 2018/1092 (Dz.Urz. UE L 170 z 12.5.2021, s.~149).
\BibItem{UAS 945}{uas945} Rozporządzenie delegowane Komisji (UE) 2019/945 z dnia 12 marca 2019 r. w sprawie bezzałogowych systemów powietrznych oraz operatorów bezzałogowych systemów powietrznych z państw trzecich (Dz.Urz. UE L 152 z 11.6.2019, s.~1).
\BibItem{UAS 947}{uas947} Rozporządzenie wykonawcze Komisji (UE) 2019/947 z dnia 24 maja 2019 r. w sprawie przepisów i procedur dotyczących eksploatacji bezzałogowych statków powietrznych (Dz.Urz. UE L 152 z 11.6.2019, s.~45).
\BibItem{FDI}{fdi} Rozporządzenie Parlamentu Europejskiego i Rady (UE) 2019/452 z dnia 19 marca 2019 r. ustanawiające ramy monitorowania bezpośrednich inwestycji zagranicznych w Unii --- uchylone z dniem 17 stycznia 2028 r. przez \Zr{fdi2026}.
\BibItem{FDI 2026}{fdi2026} Rozporządzenie Parlamentu Europejskiego i Rady (UE) 2026/1386 z dnia 17 czerwca 2026 r. w sprawie monitorowania inwestycji zagranicznych w Unii i uchylające rozporządzenie (UE) 2019/452 (Dz.Urz. UE L 2026/1386 z 26.6.2026); stosuje się od 17 stycznia 2028 r. (niektóre przepisy od 16 lipca 2026 r.); obejmuje m.in. spółki rozwijające produkty podwójnego zastosowania, wojskowe i drony.
\BibItem{TTBER}{ttber} Rozporządzenie Komisji (UE) 2026/877 z dnia 16 kwietnia 2026 r. w sprawie stosowania art.~101 ust.~3 TFUE do kategorii porozumień o transferze technologii (Dz.Urz. UE L 2026/877 z 21.4.2026); obowiązuje od 1 maja 2026 r. do 30 kwietnia 2038 r.; zastąpiło rozporządzenie (UE) nr 316/2014 --- licencje na technologię, know-how i oprogramowanie.
\BibItem{AI Act}{aiact} Rozporządzenie Parlamentu Europejskiego i Rady (UE) 2024/1689 z dnia 13 czerwca 2024 r. w sprawie ustanowienia zharmonizowanych przepisów dotyczących sztucznej inteligencji (akt w sprawie sztucznej inteligencji) (Dz.Urz. UE L 2024/1689 z 12.7.2024); systemy wyłącznie wojskowe są wyłączone (art.~2 ust.~3), systemy podwójnego zastosowania --- objęte.
\BibItem{Sankcje}{sankcje} Rozporządzenie Rady (UE) nr 833/2014 z dnia 31 lipca 2014 r. dotyczące środków ograniczających w związku z działaniami Rosji destabilizującymi sytuację na Ukrainie (Dz.Urz. UE L 229 z 31.7.2014, s.~1) oraz rozporządzenie Rady (UE) nr 269/2014 z dnia 17 marca 2014 r. (Dz.Urz. UE L 78 z 17.3.2014, s.~6) (środki ograniczające).
\BibItem{TSD}{tsd} Dyrektywa Parlamentu Europejskiego i Rady (UE) 2016/943 z dnia 8 czerwca 2016 r. w sprawie ochrony niejawnego know-how i niejawnych informacji handlowych (tajemnic przedsiębiorstwa) przed ich bezprawnym pozyskiwaniem, wykorzystywaniem i ujawnianiem (Dz.Urz. UE L 157 z 15.6.2016, s.~1).
\BibItem{WBD}{wbd} Dyrektywa Parlamentu Europejskiego i Rady (UE) 2019/1937 z dnia 23 października 2019 r. w sprawie ochrony osób zgłaszających naruszenia prawa Unii.
\end{BibList}

\section*{C. Orzecznictwo}
\begin{BibList}
\BibItem{SN 61/03}{sn6103} Uchwała składu siedmiu sędziów SN z dnia 6 listopada 2003 r., III CZP 61/03, OSNC 2004, nr 5, poz. 69 (zasada prawna) --- kara umowna należy się także wtedy, gdy wierzyciel nie poniósł szkody; dłużnik może żądać jej miarkowania (art.~484 §~2 KC).
\BibItem{SN 304/00}{sn30400} Wyrok SN z dnia 3 października 2000 r., I CKN 304/00, OSNC 2001, nr 4, poz. 59 --- wykorzystanie przez pracownika informacji, co do których przedsiębiorca nie podjął niezbędnych działań w celu zachowania ich poufności, należy traktować jako wykorzystanie powszechnej wiedzy, do której przedsiębiorca nie ma ustawowych uprawnień.
\BibItem{SN 120/07}{sn12007} Wyrok SN z dnia 8 stycznia 2008 r., II PK 120/07 --- odpowiedzialność pracownika za szkodę regulują wyczerpująco przepisy Kodeksu pracy; kara umowna (art.~483 KC) zastrzeżona w stosunku pracy nie może być dochodzona; dopuszczalna jest natomiast w umowie o zakazie konkurencji po ustaniu stosunku pracy (SN powołuje wyrok z 10 października 2003 r., I PK 528/02, OSNP 2004, nr 19, poz. 336).
\BibItem{SN 42/18}{sn4218} Uchwała SN z dnia 7 września 2018 r., III CZP 42/18 (Lex nr 2541937) --- członek zarządu nie może jako pełnomocnik reprezentować spółki z o.o. w umowie z innym członkiem zarządu (art.~210 §~1 KSH); podobnie uchwała SN z 30 stycznia 2019 r., III CZP 71/18. Skutkiem naruszenia jest bezwzględna nieważność czynności (art.~58 §~1 KC) --- tak wyrok SA w Warszawie z 9 kwietnia 2019 r., V ACa 485/18, który obie uchwały przywołuje. Przydatna przy wykładni art.~300$^{79}$ KSH.
\end{BibList}

\section*{D. Literatura}
\begin{BibList}
\BibItem{Malinowski}{malinowski} P. Malinowski (red.), \emph{Prosta spółka akcyjna. Komentarz}, Warszawa 2021.
\BibItem{Opalski}{opalski} A. Opalski, \emph{Prosta spółka akcyjna --- nowy typ spółki handlowej}, cz.~I: „Przegląd Prawa Handlowego” 2019, nr 11, s.~5--15; cz.~II: tamże, nr 12, s.~4--16.
\BibItem{ECFR}{ecfr} \emph{The Simple Joint Stock Company: Emergence of a New Close Company in Poland}, „European Company and Financial Law Review” 2024, t.~21, DOI 10.1515/ecfr-2024-0021.
\BibItem{Jara}{jara} Z. Jara (red.), \emph{Kodeks spółek handlowych. Komentarz}, C.H. Beck (wydanie aktualne).
\BibItem{Gniewek}{gniewek} E. Gniewek, P. Machnikowski (red.), \emph{Kodeks cywilny. Komentarz}, C.H. Beck (wydanie aktualne).
\BibItem{Barta}{barta} J. Barta, R. Markiewicz (red.), \emph{Ustawa o prawie autorskim i prawach pokrewnych. Komentarz}, Wolters Kluwer (wydanie aktualne).
\BibItem{Kostański}{kostanski} P. Kostański (red.), \emph{Prawo własności przemysłowej. Komentarz}, C.H. Beck (wydanie aktualne).
\BibItem{Szwaja}{szwaja} J. Szwaja (red.), \emph{Ustawa o zwalczaniu nieuczciwej konkurencji. Komentarz}, C.H. Beck (wydanie aktualne).
\end{BibList}

\section*{E. Materiały urzędowe}
\begin{BibList}
\BibItem{GUS PKD}{guspkd} Główny Urząd Statystyczny, \emph{Polska Klasyfikacja Działalności PKD 2025} (klasyfikacje.stat.gov.pl).
\BibItem{GUS CPI}{guscpi} Komunikaty Prezesa GUS w sprawie średniorocznego wskaźnika cen towarów i usług konsumpcyjnych ogółem, ogłaszane w Dzienniku Urzędowym RP „Monitor Polski” (podstawa waloryzacji kar i Cennika).
\end{BibList}

\section*{F. Dokumenty Spółki}
\begin{BibList}
\BibItem{US}{us} Umowa spółki Basilisk Systems P.S.A., wersja 3.3-S (\texttt{psa.pdf}; wersja angielska \texttt{psa\_en.pdf}).
\BibItem{UA}{ua} Umowa akcjonariuszy, wersja 1.7 (\texttt{basilisk\_sha.pdf}).
\BibItem{IP}{ip} Umowy przeniesienia praw autorskich (Umowa nr~1 i nr~2), wersja 1.6 (\texttt{basilisk\_ip.pdf}).
\BibItem{Cennik}{cennik} Cennik wynagrodzeń za rozwiązania projektowe (\texttt{basilisk\_cennik.pdf}).
\BibItem{NDA}{nda} Wzór umowy o zachowaniu poufności (\texttt{basilisk\_nda.pdf}).
\BibItem{Polityka}{polityka} Polityka bezpieczeństwa informacji (\texttt{basilisk\_polityka.pdf}).
\BibItem{Podręcznik}{podrecznik} Podręcznik podejmowania decyzji (\texttt{basilisk\_decyzje.pdf}).
\end{BibList}
 w podręczniku decyzji i komentarzu do umowy spółki.
% Odesłanie w tekście: \Zr{klucz} -> [ETYKIETA] z linkiem do pozycji.
\makeatletter
\providecommand{\BibItem}[2]{\item[\textbf{[#1]}]\def\@currentlabel{#1}\def\@currentlabelname{#1}\phantomsection\label{bib:#2}}
\makeatother
\providecommand{\Zr}[1]{[\ref{bib:#1}]}
\newenvironment{BibList}{\begin{list}{}{\setlength{\labelwidth}{2.6cm}\setlength{\leftmargin}{2.9cm}\setlength{\labelsep}{0.3cm}\setlength{\itemsep}{0.15em}\setlength{\parsep}{0pt}\small}}{\end{list}}

\chapter*{Bibliografia}
\addcontentsline{toc}{chapter}{Bibliografia}

Akty prawne podano według miejsca pierwotnej publikacji oraz ostatniego tekstu jednolitego (stan: listopad 2026), z dopiskiem „ze zm.”; przed użyciem należy sprawdzić, czy nie ogłoszono nowszego tekstu w Internetowym Systemie Aktów Prawnych (isap.sejm.gov.pl) albo w Dzienniku Ustaw (eli.gov.pl). Tezy orzeczeń podano w skrócie; przed powołaniem się na nie w sporze należy sprawdzić pełny tekst.

\section*{A. Akty prawne krajowe}
\begin{BibList}
\BibItem{KSH}{ksh} Ustawa z dnia 15 września 2000 r. --- Kodeks spółek handlowych (Dz.U. 2000 nr 94 poz. 1037 ze zm.; t.j. Dz.U. 2024 poz. 18 ze zm., w tym ustawą z dnia 23 stycznia 2026 r., Dz.U. 2026 poz. 176 --- zmiany rejestru akcjonariuszy od 18 lutego 2027 r.); przepisy o prostej spółce akcyjnej: dział IA, art.~300$^{1}$--300$^{134}$.
\BibItem{Nowela PSA}{psa2019} Ustawa z dnia 19 lipca 2019 r. o zmianie ustawy --- Kodeks spółek handlowych oraz niektórych innych ustaw (Dz.U. 2019 poz. 1655) --- wprowadzenie prostej spółki akcyjnej; ustawa w pierwotnym brzmieniu wchodziła w życie 1 marca 2020 r., a przepisy o PSA stosuje się od 1 lipca 2021 r. (po przesunięciu terminu).
\BibItem{KC}{kc} Ustawa z dnia 23 kwietnia 1964 r. --- Kodeks cywilny (Dz.U. 1964 nr 16 poz. 93; t.j. Dz.U. 2026 poz. 795 ze zm.).
\BibItem{KRO}{kro} Ustawa z dnia 25 lutego 1964 r. --- Kodeks rodzinny i opiekuńczy (Dz.U. 1964 nr 9 poz. 59; t.j. Dz.U. 2026 poz. 236 ze zm.).
\BibItem{KPC}{kpc} Ustawa z dnia 17 listopada 1964 r. --- Kodeks postępowania cywilnego (Dz.U. 1964 nr 43 poz. 296; t.j. Dz.U. 2026 poz. 468 ze zm.).
\BibItem{KP}{kp} Ustawa z dnia 26 czerwca 1974 r. --- Kodeks pracy (Dz.U. 1974 nr 24 poz. 141; t.j. Dz.U. 2025 poz. 277 ze zm.).
\BibItem{PrAut}{praut} Ustawa z dnia 4 lutego 1994 r. o prawie autorskim i prawach pokrewnych (Dz.U. 1994 nr 24 poz. 83; t.j. Dz.U. 2025 poz. 24 ze zm.).
\BibItem{PWP}{pwp} Ustawa z dnia 30 czerwca 2000 r. --- Prawo własności przemysłowej (Dz.U. 2001 nr 49 poz. 508; t.j. Dz.U. 2023 poz. 1170 ze zm.).
\BibItem{UZNK}{uznk} Ustawa z dnia 16 kwietnia 1993 r. o zwalczaniu nieuczciwej konkurencji (Dz.U. 1993 nr 47 poz. 211; t.j. Dz.U. 2026 poz. 85 ze zm.).
\BibItem{PIT}{pit} Ustawa z dnia 26 lipca 1991 r. o podatku dochodowym od osób fizycznych (Dz.U. 1991 nr 80 poz. 350; t.j. Dz.U. 2026 poz. 592 ze zm.).
\BibItem{CIT}{cit} Ustawa z dnia 15 lutego 1992 r. o podatku dochodowym od osób prawnych (Dz.U. 1992 nr 21 poz. 86 ze zm.).
\BibItem{UoR}{uor} Ustawa z dnia 29 września 1994 r. o rachunkowości (Dz.U. 1994 nr 121 poz. 591; t.j. Dz.U. 2026 poz. 522 ze zm.).
\BibItem{UKRS}{ukrs} Ustawa z dnia 20 sierpnia 1997 r. o Krajowym Rejestrze Sądowym (Dz.U. 1997 nr 121 poz. 769; t.j. Dz.U. 2025 poz. 869 ze zm.).
\BibItem{UONI}{uoni} Ustawa z dnia 5 sierpnia 2010 r. o ochronie informacji niejawnych (Dz.U. 2010 nr 182 poz. 1228; t.j. Dz.U. 2025 poz. 1209 ze zm.).
\BibItem{UOTS}{uots} Ustawa z dnia 29 listopada 2000 r. o obrocie z zagranicą towarami, technologiami i usługami o znaczeniu strategicznym dla bezpieczeństwa państwa, a także dla utrzymania międzynarodowego pokoju i bezpieczeństwa (Dz.U. 2000 nr 119 poz. 1250; t.j. Dz.U. 2023 poz. 1582 ze zm., w tym Dz.U. 2026 poz. 471 --- definicje i zezwolenia, m.in. pomoc techniczna).
\BibItem{UWT}{uwt} Ustawa z dnia 13 czerwca 2019 r. o wykonywaniu działalności gospodarczej w zakresie wytwarzania i obrotu materiałami wybuchowymi, bronią, amunicją oraz wyrobami i technologią o przeznaczeniu wojskowym lub policyjnym (Dz.U. 2019 poz. 1214; t.j. Dz.U. 2023 poz. 1743 ze zm.); koncesja na obrót technologią o przeznaczeniu wojskowym (art.~7), wymogi wobec wspólników z co najmniej 20\% akcji (art.~10 i 17)..
\BibItem{UKI}{uki} Ustawa z dnia 24 lipca 2015 r. o kontroli niektórych inwestycji (Dz.U. 2015 poz. 1272; t.j. Dz.U. 2026 poz. 47 ze zm.).
\BibItem{USank}{usank} Ustawa z dnia 13 kwietnia 2022 r. o szczególnych rozwiązaniach w zakresie przeciwdziałania wspieraniu agresji na Ukrainę oraz służących ochronie bezpieczeństwa narodowego (Dz.U. 2022 poz. 835; t.j. Dz.U. 2025 poz. 514 ze zm.).
\BibItem{USyg}{usyg} Ustawa z dnia 14 czerwca 2024 r. o ochronie sygnalistów (Dz.U. 2024 poz. 928).
\BibItem{UODO}{uodo} Ustawa z dnia 10 maja 2018 r. o ochronie danych osobowych (Dz.U. 2018 poz. 1000 ze zm.).
\BibItem{PrLot}{prlot} Ustawa z dnia 3 lipca 2002 r. --- Prawo lotnicze (Dz.U. 2002 nr 130 poz. 1112; t.j. Dz.U. 2025 poz. 1431 ze zm.).
\BibItem{PKD 2025}{pkd} Rozporządzenie Rady Ministrów z dnia 18 grudnia 2024 r. w sprawie Polskiej Klasyfikacji Działalności (PKD) (Dz.U. 2024 poz. 1936).
\BibItem{WT}{wt} Rozporządzenie Rady Ministrów z dnia 17 września 2019 r. w sprawie klasyfikacji rodzajów materiałów wybuchowych, broni, amunicji oraz wyrobów i technologii o przeznaczeniu wojskowym lub policyjnym, na których wytwarzanie lub obrót jest wymagane uzyskanie koncesji (Dz.U. 2019 poz. 1888) --- obejmuje m.in. bezzałogowe statki powietrzne i ich systemy (WT~V pkt~3).
\end{BibList}

\section*{B. Akty prawne Unii Europejskiej}
\begin{BibList}
\BibItem{RODO}{rodo} Rozporządzenie Parlamentu Europejskiego i Rady (UE) 2016/679 z dnia 27 kwietnia 2016 r. w sprawie ochrony osób fizycznych w związku z przetwarzaniem danych osobowych i w sprawie swobodnego przepływu takich danych oraz uchylenia dyrektywy 95/46/WE (ogólne rozporządzenie o ochronie danych) (Dz.Urz. UE L 119 z 4.5.2016, s.~1).
\BibItem{Dual-use}{dualuse} Rozporządzenie Parlamentu Europejskiego i Rady (UE) 2021/821 z dnia 20 maja 2021 r. ustanawiające unijny system kontroli wywozu, pośrednictwa, pomocy technicznej, tranzytu i transferu produktów podwójnego zastosowania (wersja przekształcona) (Dz.Urz. UE L 206 z 11.6.2021, s.~1, ze zm.).
\BibItem{EDF}{edf} Rozporządzenie Parlamentu Europejskiego i Rady (UE) 2021/697 z dnia 29 kwietnia 2021 r. ustanawiające Europejski Fundusz Obronny i uchylające rozporządzenie (UE) 2018/1092 (Dz.Urz. UE L 170 z 12.5.2021, s.~149).
\BibItem{UAS 945}{uas945} Rozporządzenie delegowane Komisji (UE) 2019/945 z dnia 12 marca 2019 r. w sprawie bezzałogowych systemów powietrznych oraz operatorów bezzałogowych systemów powietrznych z państw trzecich (Dz.Urz. UE L 152 z 11.6.2019, s.~1).
\BibItem{UAS 947}{uas947} Rozporządzenie wykonawcze Komisji (UE) 2019/947 z dnia 24 maja 2019 r. w sprawie przepisów i procedur dotyczących eksploatacji bezzałogowych statków powietrznych (Dz.Urz. UE L 152 z 11.6.2019, s.~45).
\BibItem{FDI}{fdi} Rozporządzenie Parlamentu Europejskiego i Rady (UE) 2019/452 z dnia 19 marca 2019 r. ustanawiające ramy monitorowania bezpośrednich inwestycji zagranicznych w Unii --- uchylone z dniem 17 stycznia 2028 r. przez \Zr{fdi2026}.
\BibItem{FDI 2026}{fdi2026} Rozporządzenie Parlamentu Europejskiego i Rady (UE) 2026/1386 z dnia 17 czerwca 2026 r. w sprawie monitorowania inwestycji zagranicznych w Unii i uchylające rozporządzenie (UE) 2019/452 (Dz.Urz. UE L 2026/1386 z 26.6.2026); stosuje się od 17 stycznia 2028 r. (niektóre przepisy od 16 lipca 2026 r.); obejmuje m.in. spółki rozwijające produkty podwójnego zastosowania, wojskowe i drony.
\BibItem{TTBER}{ttber} Rozporządzenie Komisji (UE) 2026/877 z dnia 16 kwietnia 2026 r. w sprawie stosowania art.~101 ust.~3 TFUE do kategorii porozumień o transferze technologii (Dz.Urz. UE L 2026/877 z 21.4.2026); obowiązuje od 1 maja 2026 r. do 30 kwietnia 2038 r.; zastąpiło rozporządzenie (UE) nr 316/2014 --- licencje na technologię, know-how i oprogramowanie.
\BibItem{AI Act}{aiact} Rozporządzenie Parlamentu Europejskiego i Rady (UE) 2024/1689 z dnia 13 czerwca 2024 r. w sprawie ustanowienia zharmonizowanych przepisów dotyczących sztucznej inteligencji (akt w sprawie sztucznej inteligencji) (Dz.Urz. UE L 2024/1689 z 12.7.2024); systemy wyłącznie wojskowe są wyłączone (art.~2 ust.~3), systemy podwójnego zastosowania --- objęte.
\BibItem{Sankcje}{sankcje} Rozporządzenie Rady (UE) nr 833/2014 z dnia 31 lipca 2014 r. dotyczące środków ograniczających w związku z działaniami Rosji destabilizującymi sytuację na Ukrainie (Dz.Urz. UE L 229 z 31.7.2014, s.~1) oraz rozporządzenie Rady (UE) nr 269/2014 z dnia 17 marca 2014 r. (Dz.Urz. UE L 78 z 17.3.2014, s.~6) (środki ograniczające).
\BibItem{TSD}{tsd} Dyrektywa Parlamentu Europejskiego i Rady (UE) 2016/943 z dnia 8 czerwca 2016 r. w sprawie ochrony niejawnego know-how i niejawnych informacji handlowych (tajemnic przedsiębiorstwa) przed ich bezprawnym pozyskiwaniem, wykorzystywaniem i ujawnianiem (Dz.Urz. UE L 157 z 15.6.2016, s.~1).
\BibItem{WBD}{wbd} Dyrektywa Parlamentu Europejskiego i Rady (UE) 2019/1937 z dnia 23 października 2019 r. w sprawie ochrony osób zgłaszających naruszenia prawa Unii.
\end{BibList}

\section*{C. Orzecznictwo}
\begin{BibList}
\BibItem{SN 61/03}{sn6103} Uchwała składu siedmiu sędziów SN z dnia 6 listopada 2003 r., III CZP 61/03, OSNC 2004, nr 5, poz. 69 (zasada prawna) --- kara umowna należy się także wtedy, gdy wierzyciel nie poniósł szkody; dłużnik może żądać jej miarkowania (art.~484 §~2 KC).
\BibItem{SN 304/00}{sn30400} Wyrok SN z dnia 3 października 2000 r., I CKN 304/00, OSNC 2001, nr 4, poz. 59 --- wykorzystanie przez pracownika informacji, co do których przedsiębiorca nie podjął niezbędnych działań w celu zachowania ich poufności, należy traktować jako wykorzystanie powszechnej wiedzy, do której przedsiębiorca nie ma ustawowych uprawnień.
\BibItem{SN 120/07}{sn12007} Wyrok SN z dnia 8 stycznia 2008 r., II PK 120/07 --- odpowiedzialność pracownika za szkodę regulują wyczerpująco przepisy Kodeksu pracy; kara umowna (art.~483 KC) zastrzeżona w stosunku pracy nie może być dochodzona; dopuszczalna jest natomiast w umowie o zakazie konkurencji po ustaniu stosunku pracy (SN powołuje wyrok z 10 października 2003 r., I PK 528/02, OSNP 2004, nr 19, poz. 336).
\BibItem{SN 42/18}{sn4218} Uchwała SN z dnia 7 września 2018 r., III CZP 42/18 (Lex nr 2541937) --- członek zarządu nie może jako pełnomocnik reprezentować spółki z o.o. w umowie z innym członkiem zarządu (art.~210 §~1 KSH); podobnie uchwała SN z 30 stycznia 2019 r., III CZP 71/18. Skutkiem naruszenia jest bezwzględna nieważność czynności (art.~58 §~1 KC) --- tak wyrok SA w Warszawie z 9 kwietnia 2019 r., V ACa 485/18, który obie uchwały przywołuje. Przydatna przy wykładni art.~300$^{79}$ KSH.
\end{BibList}

\section*{D. Literatura}
\begin{BibList}
\BibItem{Malinowski}{malinowski} P. Malinowski (red.), \emph{Prosta spółka akcyjna. Komentarz}, Warszawa 2021.
\BibItem{Opalski}{opalski} A. Opalski, \emph{Prosta spółka akcyjna --- nowy typ spółki handlowej}, cz.~I: „Przegląd Prawa Handlowego” 2019, nr 11, s.~5--15; cz.~II: tamże, nr 12, s.~4--16.
\BibItem{ECFR}{ecfr} \emph{The Simple Joint Stock Company: Emergence of a New Close Company in Poland}, „European Company and Financial Law Review” 2024, t.~21, DOI 10.1515/ecfr-2024-0021.
\BibItem{Jara}{jara} Z. Jara (red.), \emph{Kodeks spółek handlowych. Komentarz}, C.H. Beck (wydanie aktualne).
\BibItem{Gniewek}{gniewek} E. Gniewek, P. Machnikowski (red.), \emph{Kodeks cywilny. Komentarz}, C.H. Beck (wydanie aktualne).
\BibItem{Barta}{barta} J. Barta, R. Markiewicz (red.), \emph{Ustawa o prawie autorskim i prawach pokrewnych. Komentarz}, Wolters Kluwer (wydanie aktualne).
\BibItem{Kostański}{kostanski} P. Kostański (red.), \emph{Prawo własności przemysłowej. Komentarz}, C.H. Beck (wydanie aktualne).
\BibItem{Szwaja}{szwaja} J. Szwaja (red.), \emph{Ustawa o zwalczaniu nieuczciwej konkurencji. Komentarz}, C.H. Beck (wydanie aktualne).
\end{BibList}

\section*{E. Materiały urzędowe}
\begin{BibList}
\BibItem{GUS PKD}{guspkd} Główny Urząd Statystyczny, \emph{Polska Klasyfikacja Działalności PKD 2025} (klasyfikacje.stat.gov.pl).
\BibItem{GUS CPI}{guscpi} Komunikaty Prezesa GUS w sprawie średniorocznego wskaźnika cen towarów i usług konsumpcyjnych ogółem, ogłaszane w Dzienniku Urzędowym RP „Monitor Polski” (podstawa waloryzacji kar i Cennika).
\end{BibList}

\section*{F. Dokumenty Spółki}
\begin{BibList}
\BibItem{US}{us} Umowa spółki Basilisk Systems P.S.A., wersja 3.3-S (\texttt{psa.pdf}; wersja angielska \texttt{psa\_en.pdf}).
\BibItem{UA}{ua} Umowa akcjonariuszy, wersja 1.7 (\texttt{basilisk\_sha.pdf}).
\BibItem{IP}{ip} Umowy przeniesienia praw autorskich (Umowa nr~1 i nr~2), wersja 1.6 (\texttt{basilisk\_ip.pdf}).
\BibItem{Cennik}{cennik} Cennik wynagrodzeń za rozwiązania projektowe (\texttt{basilisk\_cennik.pdf}).
\BibItem{NDA}{nda} Wzór umowy o zachowaniu poufności (\texttt{basilisk\_nda.pdf}).
\BibItem{Polityka}{polityka} Polityka bezpieczeństwa informacji (\texttt{basilisk\_polityka.pdf}).
\BibItem{Podręcznik}{podrecznik} Podręcznik podejmowania decyzji (\texttt{basilisk\_decyzje.pdf}).
\end{BibList}
 w podręczniku decyzji i komentarzu do umowy spółki.
% Odesłanie w tekście: \Zr{klucz} -> [ETYKIETA] z linkiem do pozycji.
\makeatletter
\providecommand{\BibItem}[2]{\item[\textbf{[#1]}]\def\@currentlabel{#1}\def\@currentlabelname{#1}\phantomsection\label{bib:#2}}
\makeatother
\providecommand{\Zr}[1]{[\ref{bib:#1}]}
\newenvironment{BibList}{\begin{list}{}{\setlength{\labelwidth}{2.6cm}\setlength{\leftmargin}{2.9cm}\setlength{\labelsep}{0.3cm}\setlength{\itemsep}{0.15em}\setlength{\parsep}{0pt}\small}}{\end{list}}

\chapter*{Bibliografia}
\addcontentsline{toc}{chapter}{Bibliografia}

Akty prawne podano według miejsca pierwotnej publikacji oraz ostatniego tekstu jednolitego (stan: listopad 2026), z dopiskiem „ze zm.”; przed użyciem należy sprawdzić, czy nie ogłoszono nowszego tekstu w Internetowym Systemie Aktów Prawnych (isap.sejm.gov.pl) albo w Dzienniku Ustaw (eli.gov.pl). Tezy orzeczeń podano w skrócie; przed powołaniem się na nie w sporze należy sprawdzić pełny tekst.

\section*{A. Akty prawne krajowe}
\begin{BibList}
\BibItem{KSH}{ksh} Ustawa z dnia 15 września 2000 r. --- Kodeks spółek handlowych (Dz.U. 2000 nr 94 poz. 1037 ze zm.; t.j. Dz.U. 2024 poz. 18 ze zm., w tym ustawą z dnia 23 stycznia 2026 r., Dz.U. 2026 poz. 176 --- zmiany rejestru akcjonariuszy od 18 lutego 2027 r.); przepisy o prostej spółce akcyjnej: dział IA, art.~300$^{1}$--300$^{134}$.
\BibItem{Nowela PSA}{psa2019} Ustawa z dnia 19 lipca 2019 r. o zmianie ustawy --- Kodeks spółek handlowych oraz niektórych innych ustaw (Dz.U. 2019 poz. 1655) --- wprowadzenie prostej spółki akcyjnej; ustawa w pierwotnym brzmieniu wchodziła w życie 1 marca 2020 r., a przepisy o PSA stosuje się od 1 lipca 2021 r. (po przesunięciu terminu).
\BibItem{KC}{kc} Ustawa z dnia 23 kwietnia 1964 r. --- Kodeks cywilny (Dz.U. 1964 nr 16 poz. 93; t.j. Dz.U. 2026 poz. 795 ze zm.).
\BibItem{KRO}{kro} Ustawa z dnia 25 lutego 1964 r. --- Kodeks rodzinny i opiekuńczy (Dz.U. 1964 nr 9 poz. 59; t.j. Dz.U. 2026 poz. 236 ze zm.).
\BibItem{KPC}{kpc} Ustawa z dnia 17 listopada 1964 r. --- Kodeks postępowania cywilnego (Dz.U. 1964 nr 43 poz. 296; t.j. Dz.U. 2026 poz. 468 ze zm.).
\BibItem{KP}{kp} Ustawa z dnia 26 czerwca 1974 r. --- Kodeks pracy (Dz.U. 1974 nr 24 poz. 141; t.j. Dz.U. 2025 poz. 277 ze zm.).
\BibItem{PrAut}{praut} Ustawa z dnia 4 lutego 1994 r. o prawie autorskim i prawach pokrewnych (Dz.U. 1994 nr 24 poz. 83; t.j. Dz.U. 2025 poz. 24 ze zm.).
\BibItem{PWP}{pwp} Ustawa z dnia 30 czerwca 2000 r. --- Prawo własności przemysłowej (Dz.U. 2001 nr 49 poz. 508; t.j. Dz.U. 2023 poz. 1170 ze zm.).
\BibItem{UZNK}{uznk} Ustawa z dnia 16 kwietnia 1993 r. o zwalczaniu nieuczciwej konkurencji (Dz.U. 1993 nr 47 poz. 211; t.j. Dz.U. 2026 poz. 85 ze zm.).
\BibItem{PIT}{pit} Ustawa z dnia 26 lipca 1991 r. o podatku dochodowym od osób fizycznych (Dz.U. 1991 nr 80 poz. 350; t.j. Dz.U. 2026 poz. 592 ze zm.).
\BibItem{CIT}{cit} Ustawa z dnia 15 lutego 1992 r. o podatku dochodowym od osób prawnych (Dz.U. 1992 nr 21 poz. 86 ze zm.).
\BibItem{UoR}{uor} Ustawa z dnia 29 września 1994 r. o rachunkowości (Dz.U. 1994 nr 121 poz. 591; t.j. Dz.U. 2026 poz. 522 ze zm.).
\BibItem{UKRS}{ukrs} Ustawa z dnia 20 sierpnia 1997 r. o Krajowym Rejestrze Sądowym (Dz.U. 1997 nr 121 poz. 769; t.j. Dz.U. 2025 poz. 869 ze zm.).
\BibItem{UONI}{uoni} Ustawa z dnia 5 sierpnia 2010 r. o ochronie informacji niejawnych (Dz.U. 2010 nr 182 poz. 1228; t.j. Dz.U. 2025 poz. 1209 ze zm.).
\BibItem{UOTS}{uots} Ustawa z dnia 29 listopada 2000 r. o obrocie z zagranicą towarami, technologiami i usługami o znaczeniu strategicznym dla bezpieczeństwa państwa, a także dla utrzymania międzynarodowego pokoju i bezpieczeństwa (Dz.U. 2000 nr 119 poz. 1250; t.j. Dz.U. 2023 poz. 1582 ze zm., w tym Dz.U. 2026 poz. 471 --- definicje i zezwolenia, m.in. pomoc techniczna).
\BibItem{UWT}{uwt} Ustawa z dnia 13 czerwca 2019 r. o wykonywaniu działalności gospodarczej w zakresie wytwarzania i obrotu materiałami wybuchowymi, bronią, amunicją oraz wyrobami i technologią o przeznaczeniu wojskowym lub policyjnym (Dz.U. 2019 poz. 1214; t.j. Dz.U. 2023 poz. 1743 ze zm.); koncesja na obrót technologią o przeznaczeniu wojskowym (art.~7), wymogi wobec wspólników z co najmniej 20\% akcji (art.~10 i 17)..
\BibItem{UKI}{uki} Ustawa z dnia 24 lipca 2015 r. o kontroli niektórych inwestycji (Dz.U. 2015 poz. 1272; t.j. Dz.U. 2026 poz. 47 ze zm.).
\BibItem{USank}{usank} Ustawa z dnia 13 kwietnia 2022 r. o szczególnych rozwiązaniach w zakresie przeciwdziałania wspieraniu agresji na Ukrainę oraz służących ochronie bezpieczeństwa narodowego (Dz.U. 2022 poz. 835; t.j. Dz.U. 2025 poz. 514 ze zm.).
\BibItem{USyg}{usyg} Ustawa z dnia 14 czerwca 2024 r. o ochronie sygnalistów (Dz.U. 2024 poz. 928).
\BibItem{UODO}{uodo} Ustawa z dnia 10 maja 2018 r. o ochronie danych osobowych (Dz.U. 2018 poz. 1000 ze zm.).
\BibItem{PrLot}{prlot} Ustawa z dnia 3 lipca 2002 r. --- Prawo lotnicze (Dz.U. 2002 nr 130 poz. 1112; t.j. Dz.U. 2025 poz. 1431 ze zm.).
\BibItem{PKD 2025}{pkd} Rozporządzenie Rady Ministrów z dnia 18 grudnia 2024 r. w sprawie Polskiej Klasyfikacji Działalności (PKD) (Dz.U. 2024 poz. 1936).
\BibItem{WT}{wt} Rozporządzenie Rady Ministrów z dnia 17 września 2019 r. w sprawie klasyfikacji rodzajów materiałów wybuchowych, broni, amunicji oraz wyrobów i technologii o przeznaczeniu wojskowym lub policyjnym, na których wytwarzanie lub obrót jest wymagane uzyskanie koncesji (Dz.U. 2019 poz. 1888) --- obejmuje m.in. bezzałogowe statki powietrzne i ich systemy (WT~V pkt~3).
\end{BibList}

\section*{B. Akty prawne Unii Europejskiej}
\begin{BibList}
\BibItem{RODO}{rodo} Rozporządzenie Parlamentu Europejskiego i Rady (UE) 2016/679 z dnia 27 kwietnia 2016 r. w sprawie ochrony osób fizycznych w związku z przetwarzaniem danych osobowych i w sprawie swobodnego przepływu takich danych oraz uchylenia dyrektywy 95/46/WE (ogólne rozporządzenie o ochronie danych) (Dz.Urz. UE L 119 z 4.5.2016, s.~1).
\BibItem{Dual-use}{dualuse} Rozporządzenie Parlamentu Europejskiego i Rady (UE) 2021/821 z dnia 20 maja 2021 r. ustanawiające unijny system kontroli wywozu, pośrednictwa, pomocy technicznej, tranzytu i transferu produktów podwójnego zastosowania (wersja przekształcona) (Dz.Urz. UE L 206 z 11.6.2021, s.~1, ze zm.).
\BibItem{EDF}{edf} Rozporządzenie Parlamentu Europejskiego i Rady (UE) 2021/697 z dnia 29 kwietnia 2021 r. ustanawiające Europejski Fundusz Obronny i uchylające rozporządzenie (UE) 2018/1092 (Dz.Urz. UE L 170 z 12.5.2021, s.~149).
\BibItem{UAS 945}{uas945} Rozporządzenie delegowane Komisji (UE) 2019/945 z dnia 12 marca 2019 r. w sprawie bezzałogowych systemów powietrznych oraz operatorów bezzałogowych systemów powietrznych z państw trzecich (Dz.Urz. UE L 152 z 11.6.2019, s.~1).
\BibItem{UAS 947}{uas947} Rozporządzenie wykonawcze Komisji (UE) 2019/947 z dnia 24 maja 2019 r. w sprawie przepisów i procedur dotyczących eksploatacji bezzałogowych statków powietrznych (Dz.Urz. UE L 152 z 11.6.2019, s.~45).
\BibItem{FDI}{fdi} Rozporządzenie Parlamentu Europejskiego i Rady (UE) 2019/452 z dnia 19 marca 2019 r. ustanawiające ramy monitorowania bezpośrednich inwestycji zagranicznych w Unii --- uchylone z dniem 17 stycznia 2028 r. przez \Zr{fdi2026}.
\BibItem{FDI 2026}{fdi2026} Rozporządzenie Parlamentu Europejskiego i Rady (UE) 2026/1386 z dnia 17 czerwca 2026 r. w sprawie monitorowania inwestycji zagranicznych w Unii i uchylające rozporządzenie (UE) 2019/452 (Dz.Urz. UE L 2026/1386 z 26.6.2026); stosuje się od 17 stycznia 2028 r. (niektóre przepisy od 16 lipca 2026 r.); obejmuje m.in. spółki rozwijające produkty podwójnego zastosowania, wojskowe i drony.
\BibItem{TTBER}{ttber} Rozporządzenie Komisji (UE) 2026/877 z dnia 16 kwietnia 2026 r. w sprawie stosowania art.~101 ust.~3 TFUE do kategorii porozumień o transferze technologii (Dz.Urz. UE L 2026/877 z 21.4.2026); obowiązuje od 1 maja 2026 r. do 30 kwietnia 2038 r.; zastąpiło rozporządzenie (UE) nr 316/2014 --- licencje na technologię, know-how i oprogramowanie.
\BibItem{AI Act}{aiact} Rozporządzenie Parlamentu Europejskiego i Rady (UE) 2024/1689 z dnia 13 czerwca 2024 r. w sprawie ustanowienia zharmonizowanych przepisów dotyczących sztucznej inteligencji (akt w sprawie sztucznej inteligencji) (Dz.Urz. UE L 2024/1689 z 12.7.2024); systemy wyłącznie wojskowe są wyłączone (art.~2 ust.~3), systemy podwójnego zastosowania --- objęte.
\BibItem{Sankcje}{sankcje} Rozporządzenie Rady (UE) nr 833/2014 z dnia 31 lipca 2014 r. dotyczące środków ograniczających w związku z działaniami Rosji destabilizującymi sytuację na Ukrainie (Dz.Urz. UE L 229 z 31.7.2014, s.~1) oraz rozporządzenie Rady (UE) nr 269/2014 z dnia 17 marca 2014 r. (Dz.Urz. UE L 78 z 17.3.2014, s.~6) (środki ograniczające).
\BibItem{TSD}{tsd} Dyrektywa Parlamentu Europejskiego i Rady (UE) 2016/943 z dnia 8 czerwca 2016 r. w sprawie ochrony niejawnego know-how i niejawnych informacji handlowych (tajemnic przedsiębiorstwa) przed ich bezprawnym pozyskiwaniem, wykorzystywaniem i ujawnianiem (Dz.Urz. UE L 157 z 15.6.2016, s.~1).
\BibItem{WBD}{wbd} Dyrektywa Parlamentu Europejskiego i Rady (UE) 2019/1937 z dnia 23 października 2019 r. w sprawie ochrony osób zgłaszających naruszenia prawa Unii.
\end{BibList}

\section*{C. Orzecznictwo}
\begin{BibList}
\BibItem{SN 61/03}{sn6103} Uchwała składu siedmiu sędziów SN z dnia 6 listopada 2003 r., III CZP 61/03, OSNC 2004, nr 5, poz. 69 (zasada prawna) --- kara umowna należy się także wtedy, gdy wierzyciel nie poniósł szkody; dłużnik może żądać jej miarkowania (art.~484 §~2 KC).
\BibItem{SN 304/00}{sn30400} Wyrok SN z dnia 3 października 2000 r., I CKN 304/00, OSNC 2001, nr 4, poz. 59 --- wykorzystanie przez pracownika informacji, co do których przedsiębiorca nie podjął niezbędnych działań w celu zachowania ich poufności, należy traktować jako wykorzystanie powszechnej wiedzy, do której przedsiębiorca nie ma ustawowych uprawnień.
\BibItem{SN 120/07}{sn12007} Wyrok SN z dnia 8 stycznia 2008 r., II PK 120/07 --- odpowiedzialność pracownika za szkodę regulują wyczerpująco przepisy Kodeksu pracy; kara umowna (art.~483 KC) zastrzeżona w stosunku pracy nie może być dochodzona; dopuszczalna jest natomiast w umowie o zakazie konkurencji po ustaniu stosunku pracy (SN powołuje wyrok z 10 października 2003 r., I PK 528/02, OSNP 2004, nr 19, poz. 336).
\BibItem{SN 42/18}{sn4218} Uchwała SN z dnia 7 września 2018 r., III CZP 42/18 (Lex nr 2541937) --- członek zarządu nie może jako pełnomocnik reprezentować spółki z o.o. w umowie z innym członkiem zarządu (art.~210 §~1 KSH); podobnie uchwała SN z 30 stycznia 2019 r., III CZP 71/18. Skutkiem naruszenia jest bezwzględna nieważność czynności (art.~58 §~1 KC) --- tak wyrok SA w Warszawie z 9 kwietnia 2019 r., V ACa 485/18, który obie uchwały przywołuje. Przydatna przy wykładni art.~300$^{79}$ KSH.
\end{BibList}

\section*{D. Literatura}
\begin{BibList}
\BibItem{Malinowski}{malinowski} P. Malinowski (red.), \emph{Prosta spółka akcyjna. Komentarz}, Warszawa 2021.
\BibItem{Opalski}{opalski} A. Opalski, \emph{Prosta spółka akcyjna --- nowy typ spółki handlowej}, cz.~I: „Przegląd Prawa Handlowego” 2019, nr 11, s.~5--15; cz.~II: tamże, nr 12, s.~4--16.
\BibItem{ECFR}{ecfr} \emph{The Simple Joint Stock Company: Emergence of a New Close Company in Poland}, „European Company and Financial Law Review” 2024, t.~21, DOI 10.1515/ecfr-2024-0021.
\BibItem{Jara}{jara} Z. Jara (red.), \emph{Kodeks spółek handlowych. Komentarz}, C.H. Beck (wydanie aktualne).
\BibItem{Gniewek}{gniewek} E. Gniewek, P. Machnikowski (red.), \emph{Kodeks cywilny. Komentarz}, C.H. Beck (wydanie aktualne).
\BibItem{Barta}{barta} J. Barta, R. Markiewicz (red.), \emph{Ustawa o prawie autorskim i prawach pokrewnych. Komentarz}, Wolters Kluwer (wydanie aktualne).
\BibItem{Kostański}{kostanski} P. Kostański (red.), \emph{Prawo własności przemysłowej. Komentarz}, C.H. Beck (wydanie aktualne).
\BibItem{Szwaja}{szwaja} J. Szwaja (red.), \emph{Ustawa o zwalczaniu nieuczciwej konkurencji. Komentarz}, C.H. Beck (wydanie aktualne).
\end{BibList}

\section*{E. Materiały urzędowe}
\begin{BibList}
\BibItem{GUS PKD}{guspkd} Główny Urząd Statystyczny, \emph{Polska Klasyfikacja Działalności PKD 2025} (klasyfikacje.stat.gov.pl).
\BibItem{GUS CPI}{guscpi} Komunikaty Prezesa GUS w sprawie średniorocznego wskaźnika cen towarów i usług konsumpcyjnych ogółem, ogłaszane w Dzienniku Urzędowym RP „Monitor Polski” (podstawa waloryzacji kar i Cennika).
\end{BibList}

\section*{F. Dokumenty Spółki}
\begin{BibList}
\BibItem{US}{us} Umowa spółki Basilisk Systems P.S.A., wersja 3.3-S (\texttt{psa.pdf}; wersja angielska \texttt{psa\_en.pdf}).
\BibItem{UA}{ua} Umowa akcjonariuszy, wersja 1.7 (\texttt{basilisk\_sha.pdf}).
\BibItem{IP}{ip} Umowy przeniesienia praw autorskich (Umowa nr~1 i nr~2), wersja 1.6 (\texttt{basilisk\_ip.pdf}).
\BibItem{Cennik}{cennik} Cennik wynagrodzeń za rozwiązania projektowe (\texttt{basilisk\_cennik.pdf}).
\BibItem{NDA}{nda} Wzór umowy o zachowaniu poufności (\texttt{basilisk\_nda.pdf}).
\BibItem{Polityka}{polityka} Polityka bezpieczeństwa informacji (\texttt{basilisk\_polityka.pdf}).
\BibItem{Podręcznik}{podrecznik} Podręcznik podejmowania decyzji (\texttt{basilisk\_decyzje.pdf}).
\end{BibList}
 w podręczniku decyzji i komentarzu do umowy spółki.
% Odesłanie w tekście: \Zr{klucz} -> [ETYKIETA] z linkiem do pozycji.
\makeatletter
\providecommand{\BibItem}[2]{\item[\textbf{[#1]}]\def\@currentlabel{#1}\def\@currentlabelname{#1}\phantomsection\label{bib:#2}}
\makeatother
\providecommand{\Zr}[1]{[\ref{bib:#1}]}
\newenvironment{BibList}{\begin{list}{}{\setlength{\labelwidth}{2.6cm}\setlength{\leftmargin}{2.9cm}\setlength{\labelsep}{0.3cm}\setlength{\itemsep}{0.15em}\setlength{\parsep}{0pt}\small}}{\end{list}}

\chapter*{Bibliografia}
\addcontentsline{toc}{chapter}{Bibliografia}

Akty prawne podano według miejsca pierwotnej publikacji oraz ostatniego tekstu jednolitego (stan: listopad 2026), z dopiskiem „ze zm.”; przed użyciem należy sprawdzić, czy nie ogłoszono nowszego tekstu w Internetowym Systemie Aktów Prawnych (isap.sejm.gov.pl) albo w Dzienniku Ustaw (eli.gov.pl). Tezy orzeczeń podano w skrócie; przed powołaniem się na nie w sporze należy sprawdzić pełny tekst.

\section*{A. Akty prawne krajowe}
\begin{BibList}
\BibItem{KSH}{ksh} Ustawa z dnia 15 września 2000 r. --- Kodeks spółek handlowych (Dz.U. 2000 nr 94 poz. 1037 ze zm.; t.j. Dz.U. 2024 poz. 18 ze zm., w tym ustawą z dnia 23 stycznia 2026 r., Dz.U. 2026 poz. 176 --- zmiany rejestru akcjonariuszy od 18 lutego 2027 r.); przepisy o prostej spółce akcyjnej: dział IA, art.~300$^{1}$--300$^{134}$.
\BibItem{Nowela PSA}{psa2019} Ustawa z dnia 19 lipca 2019 r. o zmianie ustawy --- Kodeks spółek handlowych oraz niektórych innych ustaw (Dz.U. 2019 poz. 1655) --- wprowadzenie prostej spółki akcyjnej; ustawa w pierwotnym brzmieniu wchodziła w życie 1 marca 2020 r., a przepisy o PSA stosuje się od 1 lipca 2021 r. (po przesunięciu terminu).
\BibItem{KC}{kc} Ustawa z dnia 23 kwietnia 1964 r. --- Kodeks cywilny (Dz.U. 1964 nr 16 poz. 93; t.j. Dz.U. 2026 poz. 795 ze zm.).
\BibItem{KRO}{kro} Ustawa z dnia 25 lutego 1964 r. --- Kodeks rodzinny i opiekuńczy (Dz.U. 1964 nr 9 poz. 59; t.j. Dz.U. 2026 poz. 236 ze zm.).
\BibItem{KPC}{kpc} Ustawa z dnia 17 listopada 1964 r. --- Kodeks postępowania cywilnego (Dz.U. 1964 nr 43 poz. 296; t.j. Dz.U. 2026 poz. 468 ze zm.).
\BibItem{KP}{kp} Ustawa z dnia 26 czerwca 1974 r. --- Kodeks pracy (Dz.U. 1974 nr 24 poz. 141; t.j. Dz.U. 2025 poz. 277 ze zm.).
\BibItem{PrAut}{praut} Ustawa z dnia 4 lutego 1994 r. o prawie autorskim i prawach pokrewnych (Dz.U. 1994 nr 24 poz. 83; t.j. Dz.U. 2025 poz. 24 ze zm.).
\BibItem{PWP}{pwp} Ustawa z dnia 30 czerwca 2000 r. --- Prawo własności przemysłowej (Dz.U. 2001 nr 49 poz. 508; t.j. Dz.U. 2023 poz. 1170 ze zm.).
\BibItem{UZNK}{uznk} Ustawa z dnia 16 kwietnia 1993 r. o zwalczaniu nieuczciwej konkurencji (Dz.U. 1993 nr 47 poz. 211; t.j. Dz.U. 2026 poz. 85 ze zm.).
\BibItem{PIT}{pit} Ustawa z dnia 26 lipca 1991 r. o podatku dochodowym od osób fizycznych (Dz.U. 1991 nr 80 poz. 350; t.j. Dz.U. 2026 poz. 592 ze zm.).
\BibItem{CIT}{cit} Ustawa z dnia 15 lutego 1992 r. o podatku dochodowym od osób prawnych (Dz.U. 1992 nr 21 poz. 86 ze zm.).
\BibItem{UoR}{uor} Ustawa z dnia 29 września 1994 r. o rachunkowości (Dz.U. 1994 nr 121 poz. 591; t.j. Dz.U. 2026 poz. 522 ze zm.).
\BibItem{UKRS}{ukrs} Ustawa z dnia 20 sierpnia 1997 r. o Krajowym Rejestrze Sądowym (Dz.U. 1997 nr 121 poz. 769; t.j. Dz.U. 2025 poz. 869 ze zm.).
\BibItem{UONI}{uoni} Ustawa z dnia 5 sierpnia 2010 r. o ochronie informacji niejawnych (Dz.U. 2010 nr 182 poz. 1228; t.j. Dz.U. 2025 poz. 1209 ze zm.).
\BibItem{UOTS}{uots} Ustawa z dnia 29 listopada 2000 r. o obrocie z zagranicą towarami, technologiami i usługami o znaczeniu strategicznym dla bezpieczeństwa państwa, a także dla utrzymania międzynarodowego pokoju i bezpieczeństwa (Dz.U. 2000 nr 119 poz. 1250; t.j. Dz.U. 2023 poz. 1582 ze zm., w tym Dz.U. 2026 poz. 471 --- definicje i zezwolenia, m.in. pomoc techniczna).
\BibItem{UWT}{uwt} Ustawa z dnia 13 czerwca 2019 r. o wykonywaniu działalności gospodarczej w zakresie wytwarzania i obrotu materiałami wybuchowymi, bronią, amunicją oraz wyrobami i technologią o przeznaczeniu wojskowym lub policyjnym (Dz.U. 2019 poz. 1214; t.j. Dz.U. 2023 poz. 1743 ze zm.); koncesja na obrót technologią o przeznaczeniu wojskowym (art.~7), wymogi wobec wspólników z co najmniej 20\% akcji (art.~10 i 17)..
\BibItem{UKI}{uki} Ustawa z dnia 24 lipca 2015 r. o kontroli niektórych inwestycji (Dz.U. 2015 poz. 1272; t.j. Dz.U. 2026 poz. 47 ze zm.).
\BibItem{USank}{usank} Ustawa z dnia 13 kwietnia 2022 r. o szczególnych rozwiązaniach w zakresie przeciwdziałania wspieraniu agresji na Ukrainę oraz służących ochronie bezpieczeństwa narodowego (Dz.U. 2022 poz. 835; t.j. Dz.U. 2025 poz. 514 ze zm.).
\BibItem{USyg}{usyg} Ustawa z dnia 14 czerwca 2024 r. o ochronie sygnalistów (Dz.U. 2024 poz. 928).
\BibItem{UODO}{uodo} Ustawa z dnia 10 maja 2018 r. o ochronie danych osobowych (Dz.U. 2018 poz. 1000 ze zm.).
\BibItem{PrLot}{prlot} Ustawa z dnia 3 lipca 2002 r. --- Prawo lotnicze (Dz.U. 2002 nr 130 poz. 1112; t.j. Dz.U. 2025 poz. 1431 ze zm.).
\BibItem{PKD 2025}{pkd} Rozporządzenie Rady Ministrów z dnia 18 grudnia 2024 r. w sprawie Polskiej Klasyfikacji Działalności (PKD) (Dz.U. 2024 poz. 1936).
\BibItem{WT}{wt} Rozporządzenie Rady Ministrów z dnia 17 września 2019 r. w sprawie klasyfikacji rodzajów materiałów wybuchowych, broni, amunicji oraz wyrobów i technologii o przeznaczeniu wojskowym lub policyjnym, na których wytwarzanie lub obrót jest wymagane uzyskanie koncesji (Dz.U. 2019 poz. 1888) --- obejmuje m.in. bezzałogowe statki powietrzne i ich systemy (WT~V pkt~3).
\end{BibList}

\section*{B. Akty prawne Unii Europejskiej}
\begin{BibList}
\BibItem{RODO}{rodo} Rozporządzenie Parlamentu Europejskiego i Rady (UE) 2016/679 z dnia 27 kwietnia 2016 r. w sprawie ochrony osób fizycznych w związku z przetwarzaniem danych osobowych i w sprawie swobodnego przepływu takich danych oraz uchylenia dyrektywy 95/46/WE (ogólne rozporządzenie o ochronie danych) (Dz.Urz. UE L 119 z 4.5.2016, s.~1).
\BibItem{Dual-use}{dualuse} Rozporządzenie Parlamentu Europejskiego i Rady (UE) 2021/821 z dnia 20 maja 2021 r. ustanawiające unijny system kontroli wywozu, pośrednictwa, pomocy technicznej, tranzytu i transferu produktów podwójnego zastosowania (wersja przekształcona) (Dz.Urz. UE L 206 z 11.6.2021, s.~1, ze zm.).
\BibItem{EDF}{edf} Rozporządzenie Parlamentu Europejskiego i Rady (UE) 2021/697 z dnia 29 kwietnia 2021 r. ustanawiające Europejski Fundusz Obronny i uchylające rozporządzenie (UE) 2018/1092 (Dz.Urz. UE L 170 z 12.5.2021, s.~149).
\BibItem{UAS 945}{uas945} Rozporządzenie delegowane Komisji (UE) 2019/945 z dnia 12 marca 2019 r. w sprawie bezzałogowych systemów powietrznych oraz operatorów bezzałogowych systemów powietrznych z państw trzecich (Dz.Urz. UE L 152 z 11.6.2019, s.~1).
\BibItem{UAS 947}{uas947} Rozporządzenie wykonawcze Komisji (UE) 2019/947 z dnia 24 maja 2019 r. w sprawie przepisów i procedur dotyczących eksploatacji bezzałogowych statków powietrznych (Dz.Urz. UE L 152 z 11.6.2019, s.~45).
\BibItem{FDI}{fdi} Rozporządzenie Parlamentu Europejskiego i Rady (UE) 2019/452 z dnia 19 marca 2019 r. ustanawiające ramy monitorowania bezpośrednich inwestycji zagranicznych w Unii --- uchylone z dniem 17 stycznia 2028 r. przez \Zr{fdi2026}.
\BibItem{FDI 2026}{fdi2026} Rozporządzenie Parlamentu Europejskiego i Rady (UE) 2026/1386 z dnia 17 czerwca 2026 r. w sprawie monitorowania inwestycji zagranicznych w Unii i uchylające rozporządzenie (UE) 2019/452 (Dz.Urz. UE L 2026/1386 z 26.6.2026); stosuje się od 17 stycznia 2028 r. (niektóre przepisy od 16 lipca 2026 r.); obejmuje m.in. spółki rozwijające produkty podwójnego zastosowania, wojskowe i drony.
\BibItem{TTBER}{ttber} Rozporządzenie Komisji (UE) 2026/877 z dnia 16 kwietnia 2026 r. w sprawie stosowania art.~101 ust.~3 TFUE do kategorii porozumień o transferze technologii (Dz.Urz. UE L 2026/877 z 21.4.2026); obowiązuje od 1 maja 2026 r. do 30 kwietnia 2038 r.; zastąpiło rozporządzenie (UE) nr 316/2014 --- licencje na technologię, know-how i oprogramowanie.
\BibItem{AI Act}{aiact} Rozporządzenie Parlamentu Europejskiego i Rady (UE) 2024/1689 z dnia 13 czerwca 2024 r. w sprawie ustanowienia zharmonizowanych przepisów dotyczących sztucznej inteligencji (akt w sprawie sztucznej inteligencji) (Dz.Urz. UE L 2024/1689 z 12.7.2024); systemy wyłącznie wojskowe są wyłączone (art.~2 ust.~3), systemy podwójnego zastosowania --- objęte.
\BibItem{Sankcje}{sankcje} Rozporządzenie Rady (UE) nr 833/2014 z dnia 31 lipca 2014 r. dotyczące środków ograniczających w związku z działaniami Rosji destabilizującymi sytuację na Ukrainie (Dz.Urz. UE L 229 z 31.7.2014, s.~1) oraz rozporządzenie Rady (UE) nr 269/2014 z dnia 17 marca 2014 r. (Dz.Urz. UE L 78 z 17.3.2014, s.~6) (środki ograniczające).
\BibItem{TSD}{tsd} Dyrektywa Parlamentu Europejskiego i Rady (UE) 2016/943 z dnia 8 czerwca 2016 r. w sprawie ochrony niejawnego know-how i niejawnych informacji handlowych (tajemnic przedsiębiorstwa) przed ich bezprawnym pozyskiwaniem, wykorzystywaniem i ujawnianiem (Dz.Urz. UE L 157 z 15.6.2016, s.~1).
\BibItem{WBD}{wbd} Dyrektywa Parlamentu Europejskiego i Rady (UE) 2019/1937 z dnia 23 października 2019 r. w sprawie ochrony osób zgłaszających naruszenia prawa Unii.
\end{BibList}

\section*{C. Orzecznictwo}
\begin{BibList}
\BibItem{SN 61/03}{sn6103} Uchwała składu siedmiu sędziów SN z dnia 6 listopada 2003 r., III CZP 61/03, OSNC 2004, nr 5, poz. 69 (zasada prawna) --- kara umowna należy się także wtedy, gdy wierzyciel nie poniósł szkody; dłużnik może żądać jej miarkowania (art.~484 §~2 KC).
\BibItem{SN 304/00}{sn30400} Wyrok SN z dnia 3 października 2000 r., I CKN 304/00, OSNC 2001, nr 4, poz. 59 --- wykorzystanie przez pracownika informacji, co do których przedsiębiorca nie podjął niezbędnych działań w celu zachowania ich poufności, należy traktować jako wykorzystanie powszechnej wiedzy, do której przedsiębiorca nie ma ustawowych uprawnień.
\BibItem{SN 120/07}{sn12007} Wyrok SN z dnia 8 stycznia 2008 r., II PK 120/07 --- odpowiedzialność pracownika za szkodę regulują wyczerpująco przepisy Kodeksu pracy; kara umowna (art.~483 KC) zastrzeżona w stosunku pracy nie może być dochodzona; dopuszczalna jest natomiast w umowie o zakazie konkurencji po ustaniu stosunku pracy (SN powołuje wyrok z 10 października 2003 r., I PK 528/02, OSNP 2004, nr 19, poz. 336).
\BibItem{SN 42/18}{sn4218} Uchwała SN z dnia 7 września 2018 r., III CZP 42/18 (Lex nr 2541937) --- członek zarządu nie może jako pełnomocnik reprezentować spółki z o.o. w umowie z innym członkiem zarządu (art.~210 §~1 KSH); podobnie uchwała SN z 30 stycznia 2019 r., III CZP 71/18. Skutkiem naruszenia jest bezwzględna nieważność czynności (art.~58 §~1 KC) --- tak wyrok SA w Warszawie z 9 kwietnia 2019 r., V ACa 485/18, który obie uchwały przywołuje. Przydatna przy wykładni art.~300$^{79}$ KSH.
\end{BibList}

\section*{D. Literatura}
\begin{BibList}
\BibItem{Malinowski}{malinowski} P. Malinowski (red.), \emph{Prosta spółka akcyjna. Komentarz}, Warszawa 2021.
\BibItem{Opalski}{opalski} A. Opalski, \emph{Prosta spółka akcyjna --- nowy typ spółki handlowej}, cz.~I: „Przegląd Prawa Handlowego” 2019, nr 11, s.~5--15; cz.~II: tamże, nr 12, s.~4--16.
\BibItem{ECFR}{ecfr} \emph{The Simple Joint Stock Company: Emergence of a New Close Company in Poland}, „European Company and Financial Law Review” 2024, t.~21, DOI 10.1515/ecfr-2024-0021.
\BibItem{Jara}{jara} Z. Jara (red.), \emph{Kodeks spółek handlowych. Komentarz}, C.H. Beck (wydanie aktualne).
\BibItem{Gniewek}{gniewek} E. Gniewek, P. Machnikowski (red.), \emph{Kodeks cywilny. Komentarz}, C.H. Beck (wydanie aktualne).
\BibItem{Barta}{barta} J. Barta, R. Markiewicz (red.), \emph{Ustawa o prawie autorskim i prawach pokrewnych. Komentarz}, Wolters Kluwer (wydanie aktualne).
\BibItem{Kostański}{kostanski} P. Kostański (red.), \emph{Prawo własności przemysłowej. Komentarz}, C.H. Beck (wydanie aktualne).
\BibItem{Szwaja}{szwaja} J. Szwaja (red.), \emph{Ustawa o zwalczaniu nieuczciwej konkurencji. Komentarz}, C.H. Beck (wydanie aktualne).
\end{BibList}

\section*{E. Materiały urzędowe}
\begin{BibList}
\BibItem{GUS PKD}{guspkd} Główny Urząd Statystyczny, \emph{Polska Klasyfikacja Działalności PKD 2025} (klasyfikacje.stat.gov.pl).
\BibItem{GUS CPI}{guscpi} Komunikaty Prezesa GUS w sprawie średniorocznego wskaźnika cen towarów i usług konsumpcyjnych ogółem, ogłaszane w Dzienniku Urzędowym RP „Monitor Polski” (podstawa waloryzacji kar i Cennika).
\end{BibList}

\section*{F. Dokumenty Spółki}
\begin{BibList}
\BibItem{US}{us} Umowa spółki Basilisk Systems P.S.A., wersja 3.3-S (\texttt{psa.pdf}; wersja angielska \texttt{psa\_en.pdf}).
\BibItem{UA}{ua} Umowa akcjonariuszy, wersja 1.7 (\texttt{basilisk\_sha.pdf}).
\BibItem{IP}{ip} Umowy przeniesienia praw autorskich (Umowa nr~1 i nr~2), wersja 1.6 (\texttt{basilisk\_ip.pdf}).
\BibItem{Cennik}{cennik} Cennik wynagrodzeń za rozwiązania projektowe (\texttt{basilisk\_cennik.pdf}).
\BibItem{NDA}{nda} Wzór umowy o zachowaniu poufności (\texttt{basilisk\_nda.pdf}).
\BibItem{Polityka}{polityka} Polityka bezpieczeństwa informacji (\texttt{basilisk\_polityka.pdf}).
\BibItem{Podręcznik}{podrecznik} Podręcznik podejmowania decyzji (\texttt{basilisk\_decyzje.pdf}).
\end{BibList}
